We study low-rank adaptation (LoRA) in a setting small enough to measure exactly: a 2-layer, 70k-parameter character-level transformer is pretrained to sort eight digits in ascending order and then adapted to sorting in descending order and to reversal, with full fine-tuning, LoRA at ranks 1, 2, 4 and 8, last-block-only tuning and output-head-only tuning (all reimplemented), over five seeds on a CPU. With 2000 adaptation examples, LoRA at rank 4 (7,168 of the 70,016 parameters) is statistically indistinguishable from full fine-tuning on both tasks (paired $p=0.816$ and $0.374$ over five seeds), while tuning only the last block (34,496 parameters) is worse on descending sort (0.965 vs.\ 0.997 exact match) and tuning only the head fails. Accuracy on descending sort increases with rank (0.978, 0.994, 0.997, 0.997 for ranks 1, 2, 4, 8); on reversal, single seeds at ranks 1, 2 and 8 fail to optimise at the learning rate we tuned at rank 4 only and reused for all ranks, so the rank trend is confounded by optimisation instability. Retention of the pretraining task is at floor for every adapted system (sort-ascending exact match falls from 0.999 to a mean of at most 0.004, including head-only tuning and LoRA at every learning rate we tried), so the forgetting hypothesis is not supported and our design cannot discriminate between methods on it. A from-scratch model matches the pretrained one at 2000 examples; in a three-seed sweep with 100 and 300 examples, pretrained full fine-tuning is ahead of scratch, while LoRA versus full fine-tuning is inconclusive (we report no significance test for it; learning rates were tuned at 2000 examples only, and the ordering may depend on the full fine-tuning rate). Everything is small-scale and synthetic; the claims are limited to this regime.
